%% Wind-AE steady-wind solver: ported algorithm, ramp, BC/IC.
%% Author: Kwang-Il Seon
\documentclass[english,a4paper]{my_memo}
\synctex=1
\usepackage{booktabs}
\usepackage{array}
\usepackage{amsmath}
\usepackage{babel}
\emergencystretch=3em

\newcommand{\Rp}{\ensuremath{R_{\rm pl}}}
\newcommand{\Mdot}{\ensuremath{\dot M}}
\newcommand{\code}[1]{\texttt{#1}}
\newcommand{\Ftot}{\ensuremath{F_{\rm tot}}}
\newcommand{\hd}[1]{\code{\##1}}      % a #key header token

\title{The Wind-AE steady-wind solver:\\
algorithm, continuation, boundary/initial conditions,\\
and the \code{windsoln} (seed) file format}
\author{Kwang-Il Seon}
\date{2026-06-13 (updated 2026-08-19)}

\begin{document}
\maketitle

\begin{quote}\small\itshape
EXHALE includes a 1-D, steady-state, transonic Parker-wind solver: a
Fortran port of \emph{Wind-AE} (Murray-Clay et al.\ 2009; Broome et al.\
2025), under \code{src/modules/wind\_ae/}. It is used to manufacture a
warm-start initial condition for the time-dependent EXHALE run (the
\code{IC mode: windae} key, and the standalone \code{wind\_ae\_ic.x}).
This note documents what the solver computes, how it solves it (a
boundary-value relaxation from the base to the sonic point, followed by
an outward initial-value integration), how the boundary and initial
conditions are set, and how an arbitrary planet is reached from a shipped
seed by parameter continuation. It also records the structure of the
Fortran port and the gates it was validated against. Practical usage is
in the user manual (\S``Wind-AE warm-start initial conditions''); the
EXHALE-side BC/IC reference is \code{EXHALE\_BC\_and\_IC.tex}; the
code-by-code comparison with EXHALE and the other escape codes is
\code{code\_comparison.tex}; the full porting log is
\code{md\_backup/wind\_ae\_fortran\_port\_plan.md}.
\end{quote}

\tableofcontents
\bigskip

% ===================================================================== %
\section{What the solver computes}
% ===================================================================== %
Wind-AE solves the \emph{steady-state}, spherically-symmetric,
single-streamline (substellar) hydrodynamic equations for an
XUV-irradiated H/He atmosphere, from a deep base out through the sonic
point to a few planetary radii. The dependent variables are the velocity
$v$, density $\rho$, temperature $T$, the neutral fraction $Y_s$ of each
species ($s=$ H\,I, He\,I), and the photoionization column density
$N_{\rm col}$ of each species. Unlike EXHALE (a time-dependent
finite-volume code relaxed \emph{to} a steady state) Wind-AE solves
the stationary equations \emph{directly} as a two-point boundary-value
problem (BVP), so there is no time integration and no Courant limit; the
cost is that it bakes in the transonic topology (it cannot represent a
non-launching or genuinely time-dependent atmosphere).

The governing relations (code units: $r$ in \Rp, $v$ in $c_{s,0}$,
$\rho$ in $\rho_0$, $T$ in $T_0=10^4$\,K) are mass continuity, the
transonic momentum (Parker) equation, energy balance with photoionization
heating and line/recombination/Ly$\alpha$/bolometric cooling, and
photoionization--recombination equilibrium for each species:
\begin{align}
\frac{d\rho}{dr} &= -\rho\!\left(\frac{2}{r}+\frac{1}{v}\frac{dv}{dr}\right),\\
\frac{dv}{dr} &= \frac{v}{v^2-\gamma T/\mu}
   \left[\frac{2\gamma T}{\mu r}-(\gamma-1)\frac{Q}{v}-\frac{d\phi}{dr}\right],
   \label{eq:dvdr}\\
\frac{dT}{dr} &= (\gamma-1)\!\left(\frac{Q\mu}{v}+\frac{T}{\rho}\frac{d\rho}{dr}\right)
   +\frac{T}{\mu}\frac{d\mu}{dr},\\
\frac{dY_s}{dr} &= \frac{\alpha_{{\rm rec},s}\,n_e\,n_{{\rm ion},s}
   -\Gamma_s}{v\,n_{{\rm tot},s}}\frac{R_{\rm pl}}{c_{s,0}},\qquad
\frac{dN_{{\rm col},s}}{dr} = -Y_s\,n_s\,\frac{R_{\rm pl}}{N_{\rm col,0}},
\end{align}
where $Q$ is the net specific heating, $\phi$ the (optionally tidal)
potential, $\mu$ the mean molecular weight, $\gamma=5/3$,
$\Gamma_s$ the photoionization rate (primary $+$ secondary), and
$\alpha_{{\rm rec},s}$ the recombination coefficient. Equation
\eqref{eq:dvdr} is singular at the sonic point ($v^2=\gamma T/\mu$); the
solver handles it with an L'H\^opital linearization (\S\ref{sec:sonic}).

% ===================================================================== %
\section{Solution method}
% ===================================================================== %

\subsection{Relaxation on the base$\to$sonic interval (the BVP)}
The interval from the base $R_{\min}$ to the sonic point $r_s$ is mapped
to a fixed parameter $q\in[0,1]$ via $r=R_{\min}+q\,z$, where $z\equiv
r_s-R_{\min}$ is itself an unknown (a relaxed variable, constant in $r$).
On a grid of $M=1501$ points the $N_e=4+2N_{\rm spec}=8$ finite-difference
residual equations (one per variable: $v$, $z$, $\rho$, $T$, $Y_s$,
$N_{{\rm col},s}$) are solved by the Numerical Recipes \code{solvde}
Newton relaxation. Each iteration assembles the block-tridiagonal
Jacobian \code{s} by \emph{numerical} differencing of the residuals
(\code{difeq} $+$ \code{get\_eqn\_evals}, step $\delta = x/10^{7}$),
factorizes it (\code{pinvs}/\code{red}/\code{bksub}), and updates the
solution; it converges quadratically (the shipped seed relaxes to
$\|{\rm err}\|\sim10^{-11}$ in $\sim10$ iterations). The variable
ordering is permuted by \code{indexv} so the interior boundary conditions
come first, as Numerical Recipes requires.

\subsection{Outward integration (the IVP)}
Past the sonic point the problem is no longer a BVP: with $r_s$ and the
state there in hand, the same ODEs \eqref{eq:dvdr} are integrated outward
to $R_{\max}$ as an initial-value problem with an adaptive
Bulirsch--Stoer integrator (\code{odeint}\,$+$\,\code{bsstep}, with
\code{mmid}/\code{pzextr}); $\sim$1000 output points are stored. The
near-sonic singularity is past, so $dv/dr$ is evaluated directly.

\subsection{The sonic point and its linearization}\label{sec:sonic}
A transonic wind must pass smoothly through $v=c_s$, where the numerator
and denominator of \eqref{eq:dvdr} vanish together. Two boundary
conditions are imposed at $r_s$ (the last grid point): the
\emph{Bernoulli} (numerator) condition and the Mach-1 (denominator)
condition $v=c_s\cdot\code{breezeparam}$ (\code{breezeparam}$=1$ for a
true transonic wind; $<1$ gives a breeze). Approaching $r_s$, $dv/dr$ is
computed by an error-function-weighted blend of the analytic expression
and an L'H\^opital linearization
($a\,(dv/dr)^2+b\,(dv/dr)+c=0$, the positive root $=$ wind), so the
derivative stays finite and the relaxation does not stall on the $0/0$.

% ===================================================================== %
\section{Boundary conditions}
% ===================================================================== %
\paragraph{Base ($r=R_{\min}$, four conditions).} Density, temperature,
and the neutral fraction of each species are pinned to base values:
$\rho(R_{\min})=\rho_{\rm bc}$, $T(R_{\min})=T_{\rm bc}$,
$Y_s(R_{\min})=1$ (neutral base). $R_{\min}$ itself, with $\rho_{\rm bc}$
and $T_{\rm bc}$, follows from a deep, nearly-hydrostatic, bolometrically
heated base: the skin/effective temperature from $F_{\rm opt}=L_\star/(4\pi
a^2)$, a slant-path $\tau=1$ surface for $\rho$, and a microbar base
pressure (\code{base\_bcs} in the reference wrapper).
\paragraph{Sonic point ($r=r_s$, two conditions).} The Bernoulli and
Mach-1 critical-point conditions of \S\ref{sec:sonic}, plus a fixed
sonic-point column density $N_{{\rm col},s}$ per species (self-consistently
the integral of the neutral density from $r_s$ outward).

% ===================================================================== %
\section{Initial condition (the seed) and continuation}\label{sec:ramp}
% ===================================================================== %
Newton relaxation needs a nearby starting guess, and the equations admit
infinitely many breeze solutions near the transonic one; a generic guess
relaxes to a breeze, not the wind. Wind-AE therefore \emph{bootstraps}
from a single shipped converged solution and \emph{continues} it to the
target planet:
\begin{itemize}\itemsep1pt
\item \textbf{Seed.} One converged windsoln (the Murray-Clay fiducial hot
  Jupiter) is included as \code{inputdata/}\allowbreak\code{windae\_seed.csv}. It provides
  both the relaxation guess and the base BCs.
\item \textbf{Adaptive ramp (\code{ramp\_var}).} A target parameter is
  approached in multiplicative steps $p\leftarrow p(1+\delta)$; after each
  step the wind is re-relaxed from the current solution. On success the
  step accepts and $\delta$ doubles; on failure the guess is restored,
  the parameter reset, and $\delta$ halves. The run aborts a parameter
  only after many consecutive failures.
\item \textbf{Ramp order (\code{ramp\_to}).} $F_{\rm XUV}\to$ gravity
  ($R_{\rm pl}$, $M_{\rm pl}$) $\to$ star ($M_\star$, $a$, $L_\star$),
  each by \code{ramp\_var}. The converged solution of one step is the
  guess for the next (in the port this is an in-memory hand-off; the
  reference wrapper round-trips it through \code{guess.inp}).
\end{itemize}
During the ramp the base boundary conditions can either be held fixed at
the seed values (\emph{static-BC mode}) or re-converged at every step
(\emph{self-consistent-BC mode}); the port implements both, selected by an optional
\code{static\_bcs} flag on \code{ramp\_var}/\code{ramp\_to} (absent or
\code{.true.}\ $=$ static-BC mode; \code{.false.}\ $=$ self-consistent-BC mode).

\subsection{Self-consistent base BCs}\label{sec:c2}
The static-BC ramp suffices for hot Jupiters near the seed, but as the
ramp moves toward a strongly-bound, highly-irradiated planet the seed's
base density, pressure, and sonic-point column density become physically
inconsistent with the new gravity and flux, and the relaxation stalls.
The self-consistent-BC ramp restores consistency by recomputing the lower boundary as the
parameters move. It has three pieces, ported from the reference wrapper.

\paragraph{Closed-form base state (\code{base\_bcs}).} Taking \Rp\ as the
optical slant-path $\tau=1$ surface, with a bolometrically heated,
near-isothermal layer below the wind, the base temperature and density
follow in closed form from the stellar optical flux
$F_{\rm opt}=L_\star/(4\pi a^2)$:
\begin{align}
T_{\rm skin} &= \left[\frac{F_{\rm opt}\,(\kappa_{\rm opt}+\kappa_{\rm IR}/4)}
   {2\,\sigma_{\rm SB}\,\kappa_{\rm IR}}\right]^{1/4}, \qquad
T_{\rm eff} = \left[\frac{F_{\rm opt}}{4\,\sigma_{\rm SB}}\right]^{1/4},\\
\rho_{R_{\rm pl}} &= \frac{1}{\kappa_{\rm opt}}
   \sqrt{\frac{\mu\,G\,M_{\rm pl}}{8\,R_{\rm pl}^3\,k_B\,T_{\rm skin}}},
\end{align}
with default opacities $\kappa_{\rm opt}=4\times10^{-3}$,
$\kappa_{\rm IR}=10^{-2}\,{\rm cm^2\,g^{-1}}$. The simulation base is set at
the vertical IR $\tau=1$ radius $R_{\rm IR}$ if that lies above \Rp,
otherwise at the radius $R_{\rm mbar}$ where the hydrostatic pressure
reaches a microbar reference; $\rho$ and $T$ there are returned in code
units. The base mean molecular weight blends the molecular value below the
wind with the atomic value inside it,
\[
\mu = \frac{m_H}{D}\,(1-b) + b\,\mu_{\rm mol}, \qquad
D = \sum_s \frac{m_H}{m_s}\,X_s\,(2-Y_s), \qquad \mu_{\rm mol}=2.3\,m_H,
\]
where $b$ is the bolometric-transition flag ($1$ on, $0$ off) and $X_s$ the
mass fraction. This routine is \textbf{bit-identical to the reference} on
the seed and on HD\,209458b/HD\,189733b: i.e.\ on both the $b=1$
(molecular, $\mu=2.3\,m_H$) and $b=0$ (atomic, $\mu=m_H/D$) branches.

\paragraph{Sonic-point column density (\code{self\_consistent\_Ncol},
\code{converge\_Ncol\_sp}).} The column-density BC is made self-consistent
by integrating the current solution outward and setting
\[
N_{{\rm col},s} = \frac{1}{N_{\rm col,0}}\int_{r_s}^{R_{\max}}
   n_{{\rm neut},s}(r)\,dr,
\]
then re-relaxing; the set\,$\to$\,solve\,$\to$\,recompute loop iterates
until the mean relative change is below $8\%$. The ported goal matches the
reference to $\sim$4\% (within that self-consistency tolerance: the
small gap is the outward-integration extent, not a physics error).

\paragraph{Applying and ramping (\code{ramp\_base\_bcs}).} The
\code{base\_bcs} target is first attempted as a single jump; if the
relaxation fails it is reached by ramping $R_{\min}$, $T_{\rm bc}$, and
$\rho_{\rm bc}$ individually with adaptive sub-steppers (\code{ramp\_bc},
the boundary-value analogue of \code{ramp\_var}).

\paragraph{Molecular/atomic transition (\code{converge\_mol\_atomic}).} As
the ramp moves to a strongly-bound planet the base sinks \emph{into} the
wind. The transition finder (a port of the reference \code{erf\_velocity})
locates the first grid point where photoionization heating exceeds the
magnitude of PdV (adiabatic-expansion) cooling: where the wind launches;
$\code{heat}=\sum_s n_{{\rm neut},s}\,h_s$ (the rates for each species from
\code{wae\_glq\_rates\_eval}) and
$\code{cool}=(k_B T/\mu_{\rm mol}m_H)\,v\,d\rho/dr$. When that point reaches
the base (within $\sim$10 cells), \code{turn\_off\_bolo} ramps
\code{bolo\_heat\_cool}$:1\to0$; because the molecular smoothing erf is
scaled by that flag, turning it off also drops $\mu$ to its atomic value:
i.e.\ the $b=1\to b=0$ transition the target planet requires.

\paragraph{Orchestration.} With \code{static\_bcs=.false.},
\code{ramp\_var} runs the \emph{full} base-BC re-convergence
(\code{ramp\_base\_bcs} $+$ \code{converge\_mol\_atomic} $+$
\code{converge\_Ncol\_sp})
\emph{reactively}: only when a step stalls (five consecutive failures);
and the cheap molecular/atomic check \code{converge\_mol\_atomic} alone runs
\emph{proactively} (every ten successful steps: it does no relaxation
unless it actually turns the bolometric layer off, so a seed-adjacent planet
pays almost nothing). \code{ramp\_to} does a final full pass at the end.
This split keeps seed-adjacent planets fast (HD\,209458b ramps in
$\sim$40\,s) while turning the molecular layer off early enough for
far-from-seed planets (HD\,189733b in $\sim$4\,min). The reference wrapper's
expensive success-path re-convergence (\code{converge\_bcs}) likewise
defaults off.

\paragraph{Status (2026-06-13): the self-consistent-BC ramp converges HD\,189733b.} On
HD\,189733b, far from the seed ($F_{\rm XUV}$ rises $\sim$22$\times$, the
surface gravity roughly doubles), the base-BC re-convergence fires as the
ramp proceeds, and on the $M_{\rm pl}$ ramp the transition reaches the base
so \code{converge\_mol\_atomic} turns the bolometric layer off
($\code{bolo\_heat\_cool}:1\to0$, matching HD\,189733b's atomic solution).
The relaxation then tightens to $\|{\rm err}\|\sim9\times10^{-11}$ and the
full six-parameter ramp completes with a physical wind ($v>0$ throughout);
the in-process bridge runs the same ramp and writes a valid HD\,189733b IC.
\emph{Caveat:} EXHALE's subsequent time-integration of HD\,189733b is
limited by a separate, pre-existing base-breathing instability near
$r\sim1.07\,$\Rp\ (unrelated to the Wind-AE IC), so a fully converged EXHALE
warm start for that planet still depends on resolving that EXHALE-side issue.

% ===================================================================== %
\section{The Fortran port}
% ===================================================================== %
The reference Wind-AE solver core is C ($\sim$6.3k lines incl.\ Numerical
Recipes) driven by a Python wrapper that does the continuation. EXHALE's
port reimplements both in Fortran, self-contained under
\code{src/modules/wind\_ae/} (the only file that touches EXHALE's globals
is the bridge of \S\ref{sec:use}). Module map:

\begin{center}\small
\begin{tabular}{@{}lp{8.4cm}@{}}
\toprule
Module(s) & Role \\
\midrule
\code{wae\_config,\_types,\_grid,\_status} & sizes/constants, derived
  types, the shared $q$-grid, the graceful-failure flag \\
\code{wae\_rate\_coeffs,\_rnew,\_rrec,\_fe,\_ion\_pots,\_lc\_*} &
  auto-generated atomic-data tables (secondary-ionization, recombination,
  ionization potentials, line-cooling coefficients) \\
\code{wae\_spectrum} & GLQ multifrequency spectrum + planet parameters \\
\code{wae\_glq\_rates} & photoionization/heating rates (incl.\ the
  column-density rate cache) \\
\code{wae\_soe} & physics: $\mu$, $\gamma$, recombination, $Q$, all the
  $d/dr$'s, and the sonic-point linearization \\
\code{wae\_eqns} & residual equations $+$ the \code{eval\_eqn} dispatcher \\
\code{wae\_difeq} & numerical-Jacobian assembly of the \code{solvde} matrix \\
\code{wae\_nr} & the relaxation stack \code{solvde}/\code{pinvs}/\code{red}/\code{bksub} \\
\code{wae\_intode} & outward Bulirsch--Stoer integration \\
\code{wae\_params,\_continuation} & input reading; seed load,
  \code{ramp\_var}/\code{ramp\_to}, and the self-consistent base-BC handling
  (\code{base\_bcs}, \code{ramp\_base\_bcs}, \code{self\_consistent\_Ncol},
  \code{converge\_Ncol\_sp}) \\
\code{wae\_ic\_writer} & project the solution onto the EXHALE grid, write \code{*\_IC.txt} \\
\code{wae\_exhale\_input} & map an EXHALE \code{input.inp} to Wind-AE parameters \\
\bottomrule
\end{tabular}
\end{center}

\paragraph{Validation (against the C/Python reference).} Each layer was
gated with a trace-driver comparison:
the photoionization/heating rates and each species's $dY_s/dr$ are
\emph{bit-identical}; the heating/cooling $Q$ and the residual equations
match to machine precision except where the libm \code{erf} differs from
gfortran's by $\sim$1 ulp (a benign $\sim10^{-10}$, washed out at the
Newton fixed point); the \textbf{full windsoln} (relaxation $+$ outward
integration) reproduces the C solution to relative
$\le2\times10^{-11}$ on three planets; and the continuation ramp
reproduces the Python wrapper's mass-loss rate to five digits. A subtle
quirk of the reference (the rate cache keys on column density only and
ignores $Y_s$ changes) had to be reproduced exactly for the numerical
Jacobian to match.

\paragraph{The bit-exact gates need one line restored (2026-08-19).} Those
comparisons were run with Wind-AE's own \code{defs.h} constants. The port now
uses the CODATA values instead, so that the IC it hands over is built on the
same constants as the rest of EXHALE
($G=6.67430\times10^{-8}$, $k_B=1.380649\times10^{-16}$,
$m_{\rm H}=1.67353284\times10^{-24}$; Update\_EXHALE \S63.1). The differences
are $2.6\times10^{-4}$, $6.5\times10^{-6}$ and $1.4\times10^{-4}$ relative, so
agreement with the reference is now at that level rather than bit-exact.
\code{wae\_params.f90} carries the original three values in a comment; restore
them to reproduce the gates above verbatim.

% ===================================================================== %
\section{Using it in EXHALE}\label{sec:use}
% ===================================================================== %
\paragraph{In-process: \code{IC mode: windae}.} With this key in
\code{input.inp} (and \code{Load IC?\ False}), \code{init.f90} calls
\code{wae\_generate\_ic} (module \code{wae\_exhale\_bridge}), which maps
the planet's EXHALE parameters to Wind-AE
($F_{\rm XUV}=J_{\rm XUV}$, $L_\star$ from \code{Stellar Teff/radius},
$R_{\rm pl}=R_0$, the He/H ratio, the base density $n_0$ and temperature
$T_0$, and \code{Domain mode}$\to$tidal on/off), ramps the seed to the
planet, solves the wind, and writes the IC onto the EXHALE grid
(\code{output/*\_IC.txt}); \code{load\_IC} then ingests it. The whole
input\,$\to$\,solve\,$\to$\,IC\,$\to$\,run is one \code{EXHALE.x}
invocation. Because EXHALE re-solves the ionization (and metals) from the
first step, the Wind-AE IC is only a \emph{warm start}: finish with
\code{Solver: Newton} for a quantitative \Mdot.
\paragraph{Standalone: \code{wind\_ae\_ic.x}.} \code{make wind\_ae\_ic}
builds a separate executable that does the same out of process, reading an
\code{input.inp} and writing the IC files on a grid taken from an
\code{IC\_dump.txt}; run EXHALE on the result with \code{Load IC?\ True}.
The standalone tool shares the solver modules but not the EXHALE bridge,
so it stays self-contained.

% ===================================================================== %
\section{Scope and limitations}\label{sec:limits}
% ===================================================================== %
\begin{itemize}\itemsep1pt
\item \textbf{H/He only.} The port fixes $N_{\rm spec}=2$ (H\,I, He\,I):
  one ionization stage per species. It does \emph{not} model the trace
  metals, Fe/Ca cooling, the He triplet, or the transit spectrum; those
  remain EXHALE's job. For an IC this is sufficient: EXHALE adds the
  metals in equilibrium from the first step, and the IC writer rescales
  the density to EXHALE's exact He/H.
\item \textbf{Continuation BCs.} By default (\code{static\_bcs=.false.})
  the in-process ramp uses the self-consistent-BC ramp (\S\ref{sec:c2}):
  \code{base\_bcs} (bit-exact),
  \code{converge\_Ncol\_sp} (within tolerance), \code{converge\_mol\_atomic}
  (bolometric-layer turn-off), and the reactive/proactive re-convergence in
  \code{ramp\_var}/\code{ramp\_to}. This converges seed-adjacent hot Jupiters
  \emph{and} strongly-bound, far-from-seed planets such as HD\,189733b (which
  stalled under the older static-BC ramp). The static-BC path remains
  available as \code{static\_bcs=.true.}. Note this concerns the
  \emph{Wind-AE} ramp only; whether EXHALE then time-integrates a given
  planet cleanly is a separate matter (e.g.\ HD\,189733b hits an EXHALE-side
  base-breathing instability regardless of the IC source).
\item \textbf{Substellar geometry, Coriolis cutoff.} Wind-AE is a single
  substellar streamline integrated only to a few \Rp; beyond the Coriolis
  turning radius 3-D effects dominate. The mass-loss rate carries the
  Murray-Clay $0.33$ geometric factor when compared to a global rate.
\item \textbf{erf $\sim$1 ulp.} The library error function differs from
  gfortran's at the last bit; irrelevant for an IC, noted for the record.
\end{itemize}


% ===================================================================== %
% ------ The windsoln (seed) file format (merged from the former ------ %
% ------ windae_data_format.tex, same date)                      ------ %
% ===================================================================== %
\clearpage
\section*{The \code{windsoln} (seed) file format}
\addcontentsline{toc}{section}{The \code{windsoln} (seed) file format}

\begin{quote}\small\itshape
A converged Wind-AE wind solution is stored as a single CSV-like file: a
block of comment-prefixed header lines that carry every parameter,
boundary condition, and the stellar spectrum, followed by the solution
profile itself as a table of normalized (code-unit) columns. The same
format is used by (i) the original Wind-AE Python wrapper (Murray-Clay et
al.\ 2009; Broome et al.\ 2025), (ii) the included solution grid under
\code{inputdata/windae\_grid/}, and (iii) EXHALE's Fortran port, whose
\code{load\_seed} reads such a file as the starting point of a
continuation and whose \code{dump\_seed} writes one back. This note is the
authoritative field-by-field reference, taken from the Python reader
(\code{wrapper/wrapper\_utils/windsoln.py}).
\end{quote}

\section{File layout}
A windsoln file has three kinds of lines, in order:
\begin{enumerate}\itemsep2pt
  \item \textbf{Header lines} \code{\#key: v1,v2,\dots}, one per key,
        carrying scalars, parameters, BCs, flags, and spectrum metadata
        (Table~\ref{tab:hdr}).
  \item \textbf{Spectrum rows} \code{\#\# e,wPhi,$\sigma_1$,$\sigma_2$}:
        \code{npts} of them, holding the stellar spectrum sampled at
        the Gauss--Legendre quadrature points (\S\ref{sec:spec}).
  \item \textbf{Data rows}: the solution profile, one grid point per
        line, columns named by \hd{vars} and normalized by \hd{scales}
        (\S\ref{sec:grid}).
\end{enumerate}
Any other line beginning with \code{\#} is a free comment and is skipped.
A minimal excerpt:
\begin{verbatim}
#nspecies: 2
#vars: r,rho,v,T,Ys_HI,Ys_HeI,Ncol_HI,Ncol_HeI,q,z
#scales: 8.34e9,1.0e-15,9.08e5,1.0e4,1,1,1.0e17,1.0e17,1,8.34e9
#plnt_prms: 2.348e30,8.34e9,1.870e33,4.937e11,2.42e4,3.828e33
#phys_prms: 0.800,0.200,HI,HeI,1.673e-24,6.646e-24,2.30
#bcs: 1.030,7.755,2.466e4,0.198,1.0,1.0,1.08e-3,2.05e-4,1.85e-5,1.50e-5
#tech: 1.0,1.0e2,2.5,0.99
#flags: 1,1.00000,0.00000,1
#add_prms: 0
#spec_prms: 59,2,2009-01-01,scaled-solar,full,5.99e-8,9.12e-6, ... ,2.18e-11,3.94e-11
## 2.150e-9,1.771e6,4.258e-24,1.267e-22
## ... (59 rows) ...
1.030,2.466e4,2.232e-6,0.198,1.0,1.0,7.34e3,4.70e2,0.0,2.76
... (one row per grid point) ...
\end{verbatim}

\section{Header reference}
Table~\ref{tab:hdr} lists every header key. $N$ is the number of species
(\hd{nspecies}; $N=2$ for the included H/He data). Planet parameters are in
cgs; boundary conditions and the data grid are in the code units of
Table~\ref{tab:grid}.

\begin{table}[t]\centering\small
\caption{Header lines of a windsoln file. Field order is exact;
``$\times N$'' marks a block repeated for each species.}
\label{tab:hdr}
\begin{tabular}{@{}l l p{8.4cm}@{}}
\toprule
Key & Fields & Meaning (and units) \\
\midrule
\hd{nspecies}  & $N$ & number of species ($N=2$: H\,\textsc{i}, He\,\textsc{i}). \\
\hd{vars}      & names & data-grid column labels (Table~\ref{tab:grid}). \\
\hd{scales}    & $N_{\rm col}$ & cgs value of one code unit for each data column. \\
\hd{plnt\_prms}& 6 & $M_{\rm p}$[g], $\Rp$[cm], $M_\star$[g], $a$[cm],
                     $\boldsymbol{\Ftot}$\,[erg\,cm$^{-2}$s$^{-1}$],
                     $L_\star$[erg\,s$^{-1}$]. \emph{Field~5 is the total
                     ionizing flux, not $T_{\rm eq}$.} \\
\hd{phys\_prms}& $3N{+}1$ & mass fractions $X_s$ ($\times N$); species names
                     ($\times N$); atomic masses [g] ($\times N$);
                     then \code{molec\_adjust} (mean-molecular-weight boost
                     of the molecular base layer, $\approx2.3$). \\
\hd{bcs}       & $2N{+}6$ & $R_{\rm min},R_{\rm max}$ [$\Rp$];
                     $\rho_{\rm min}$ [$\rho_0$]; $T_{\rm min}$ [$T_0$];
                     base neutral fractions $Y_{s,\rm min}$ ($\times N$);
                     sonic-point columns $N_{{\rm col},s}$ [$N_{\rm col,0}$]
                     ($\times N$); \code{erf\_drop} (2: the erf-velocity
                     molecular$\to$atomic transition). \\
\hd{tech}      & 4 & \code{breezeparam}, \code{rapidity}, \code{erfn},
                     \code{mach\_limit} (relaxation/solver controls). \\
\hd{flags}     & 4 & \code{lyacool} (0/1); \textbf{\code{tidalforce}}
                     (0 = spherical, 1 = tidal/Roche); \textbf{\code{bolo\_heat\_cool}}
                     (0 = atomic base, 1 = molecular bolometric layer);
                     \code{integrate\_outward} (0/1). \\
\hd{add\_prms} & $1{+}2n$ & $n$, then $n$ $(\text{name},\text{value})$ pairs
                     of optional extra parameters ($n=0$: none). \\
\hd{spec\_prms}& meta & \code{npts}, $N$, date, source file
                     (\code{scaled-solar}), kind (\code{full}),
                     window[2], resolved[2], normalized[2] (wavelength
                     bounds [cm]), ionization potentials $\chi_s$ [erg]
                     ($\times N$). See \S\ref{sec:spec}. \\
\bottomrule
\end{tabular}
\end{table}

\section{The embedded stellar spectrum}
\label{sec:spec}
The stellar XUV spectrum that the solution was computed with is stored in
the file itself, in two parts:
\begin{itemize}\itemsep2pt
  \item \hd{spec\_prms}: metadata: the number of quadrature points
        \code{npts}, the source label (here \code{scaled-solar}, a
        solar-like SED), the spectral kind (\code{full}, i.e.\ resolved
        rather than monochromatic), the integration window in wavelength
        ($\approx[6,912]\,$\AA, i.e.\ X-ray to the H Lyman limit), and the
        the ionization potentials $\chi_s$ for each species (e.g.\
        $2.18\times10^{-11}\,$erg $=13.6$\,eV for H,
        $3.94\times10^{-11}\,$erg $=24.6$\,eV for He).
  \item the \code{\#\#} rows: the spectrum proper, one row per
        quadrature point: photon energy $hc/\lambda_i$ [erg]; the
        \emph{normalized} weight $w_i\Phi_{\lambda_i}/\Ftot$; and the
        photoionization cross sections $\sigma_{\lambda_i,s}$ [cm$^2$] for
        each species.
\end{itemize}
The crucial point is the separation of \emph{shape} and \emph{amplitude}:
the \code{\#\#} weights are normalized by $\Ftot$, so they encode only the
spectral \emph{shape} (which photon energies carry what fraction of the
ionizing flux), while the absolute level is set by the single scalar
$\Ftot$ in \hd{plnt\_prms}. This is why the included grid can keep one fixed
\code{scaled-solar} shape across all $\sim$1000 solutions and vary only
$\Ftot$ (and $M_{\rm p},\Rp$) from planet to planet. Verified: \hd{spec\_prms}
is byte-identical across the entire grid.

\medskip\noindent\textbf{Caveat (spectrum dependence).} The wind structure
depends on the SED shape, not only on $\Ftot$. The included data assume a
\code{scaled-solar} spectrum; an EXHALE run that instead assumes, say, a
power law produces different heating/ionization profiles. A Wind-AE
warm-start is therefore only self-consistent with an EXHALE run that uses
the \emph{same} spectrum. Conversely, a power-law EXHALE model failing to
converge does \emph{not} imply that no steady wind exists for that planet:
only that that particular spectrum/numerics combination did not settle.

\section{The solution grid}
\label{sec:grid}
Each data row is one grid point; the columns are named by \hd{vars} and are
stored in dimensionless code units. The physical (cgs) value is recovered
as $\text{(physical)} = \text{(code)}\times\text{(scale)}$, with the scale
taken from \hd{scales} (Table~\ref{tab:grid}). The independent variable is
the relaxation coordinate $q\in[0,1]$ mapping the base to the sonic point;
the physical radius is $r = R_{\rm min} + q\,z$, where $z$ (constant along
the solution) is the base$\to$sonic span in $\Rp$.

\begin{table}[t]\centering\small
\caption{Data-grid columns (\hd{vars}) and their code$\to$cgs scales
(\hd{scales}). $\rho_0{=}10^{-15}$, $c_{s,0}{=}\sqrt{k_BT_0/m_{\rm H}}
{=}9.08\times10^{5}$, $T_0{=}10^{4}$, $N_{\rm col,0}{=}10^{17}$.}
\label{tab:grid}
\begin{tabular}{@{}l l l p{6.0cm}@{}}
\toprule
Column & Symbol & Scale & Meaning \\
\midrule
\code{r}        & $r$            & $\Rp$              & radius \\
\code{rho}      & $\rho$         & $\rho_0$           & mass density [g\,cm$^{-3}$] \\
\code{v}        & $v$            & $c_{s,0}$          & radial velocity [cm\,s$^{-1}$] \\
\code{T}        & $T$            & $T_0$              & temperature [K] \\
\code{Ys\_HI}   & $Y_{\rm HI}$   & 1                  & H neutral fraction \\
\code{Ys\_HeI}  & $Y_{\rm HeI}$  & 1                  & He neutral fraction \\
\code{Ncol\_HI} & $N_{\rm HI}$   & $N_{\rm col,0}$    & H column to the star [cm$^{-2}$] \\
\code{Ncol\_HeI}& $N_{\rm HeI}$  & $N_{\rm col,0}$    & He column to the star [cm$^{-2}$] \\
\code{q}        & $q$            & 1                  & relaxation coordinate $\in[0,1]$ \\
\code{z}        & $z$            & $\Rp$              & base$\to$sonic span ($r=R_{\rm min}+qz$) \\
\bottomrule
\end{tabular}
\end{table}

\section{Use in the EXHALE port}
\code{src/modules/wind\_ae/wae\_continuation.f90} reads and writes this
format:
\begin{itemize}\itemsep2pt
  \item \code{load\_seed} parses \hd{plnt\_prms}, \hd{phys\_prms},
        \hd{bcs}, \hd{tech}, \hd{flags}, and the first $m$ data rows (into
        the relaxation arrays). It \emph{ignores} \hd{scales},
        \hd{add\_prms}, \hd{spec\_prms}, and the \code{\#\#} rows: the
        spectrum is loaded separately from
        \code{inputdata/windae\_spectrum.inp}, which is the same
        \code{scaled-solar} data in the equivalent
        ``\code{\# KEY: \dots}'' header form.
  \item \code{dump\_seed} is the exact inverse: it reconstructs the full
        header (including \hd{scales}, \hd{spec\_prms}, and the
        \code{\#\#} block from the in-memory spectrum arrays) and writes
        the $m$ relaxation rows, so a converged solution can be banked as a
        reusable seed.
  \item \textbf{Seed library and picker.} \code{inputdata/windae\_grid/}
        holds the Broome et al.\ (2025) grid ($\sim$1000 solutions spanning
        $M_{\rm p}\!\in\![3,528]\,M_\oplus$ and
        $\Rp\!\in\![1.9,22.6]\,R_\oplus$, all with \code{tidalforce}=1 and
        the fixed \code{scaled-solar} spectrum), plus the eight
        \code{starting\_solutions}. With \code{Wind-AE seed: grid} in
        \code{input.inp}, \code{pick\_nearest\_seed} reads the grid manifest
        (one row of \code{path}, $M_{\rm p}$, $\Rp$, $\Ftot$ per file;
        rebuilt by \code{make\_windae\_manifest.sh}) and selects the
        solution closest in $(\log M_{\rm p},\log\Rp)$, so the continuation
        starts from a near neighbor rather than the single distant fiducial
        seed. \code{Wind-AE seed out:} dumps the converged solution back
        into this format.
\end{itemize}

\end{document}
